%-------------------------------------------------------------------------------
% Document: OpenCTM Format Specification
% Author:   Marcus Geelnard
% Compile:  pdflatex FormatSpecification.tex
%-------------------------------------------------------------------------------
% Note: You need a LaTeX environment to build this document as PDF. The
% recommended way is to install TeX Live (http://www.tug.org/texlive/) and
% a decent LaTeX editor (e.g. texmaker, LEd, etc).
%
% Ubuntu:   sudo apt-get install texlive-full
% Mac OS X: http://www.tug.org/mactex/  (MacTeX.mpkg.zip)
%
% To build the PDF document, run pdflatex twice on this .tex file (in order to
% correctly build the TOC).
%-------------------------------------------------------------------------------

% Use the OpenCTM TeX style
% Document class
\documentclass[11pt,a4paper]{report}

% Packages
\usepackage[latin1]{inputenc}
\usepackage{amsmath}
\usepackage{amsfonts}
\usepackage{amssymb}
\usepackage{listings}
\usepackage{color}
\usepackage{times}

% Paragraph styles
\raggedright
\usepackage{parskip}

% Code listings
\lstset{frame=single,frameround=tttt,backgroundcolor=\color{code},%
        language=C,basicstyle={\ttfamily\small},%
        breaklines,breakindent=0pt,postbreak=\space\space\space\space}

% Colors
\definecolor{link}{rgb}{0.6,0.0,0.0}
\definecolor{code}{rgb}{0.9,0.9,1.0}
\definecolor{codeA}{rgb}{0.9,1.0,0.9}
\definecolor{codeB}{rgb}{1.0,0.9,0.9}

% PDF specific document properties
% hyperref (bookmarks, links etc) - use this package last
\usepackage[colorlinks=true,linkcolor=link,bookmarks=true,bookmarksopen=true,%
            pdfhighlight=/N,bookmarksnumbered=true,bookmarksopenlevel=1,%
            pdfview=FitH,pdfstartview=FitH]{hyperref}
\hypersetup{pdftitle={OpenCTM Developers Manual}}
\hypersetup{pdfauthor={Marcus Geelnard}}
\hypersetup{pdfkeywords={OpenCTM,manual}}


% Document properties
\author{Marcus Geelnard}
\title{OpenCTM Format Specification}

% Document contents
\begin{document}

%--[ Title page ]---------------------------------------------------------------

\begin{titlepage}
\begin{center}
~
\vspace{5cm}

\includegraphics[width=10.0cm]{logo.pdf}
\vspace{0.4cm}

{\large File format version 5}

\vspace{1.0cm}

{\Large Format Specification}
\vspace{1.5cm}

Copyright \copyright \ 2009 Marcus Geelnard
\end{center}
\end{titlepage}


%--[ Table of contents ]--------------------------------------------------------

\tableofcontents


%-------------------------------------------------------------------------------

\chapter{Overview}
This document describes version 5 of the OpenCTM file format.

\section{Data formats}
All integer fields are stored in 32-bit little endian format (least significant
byte first).

All floating point fields are stored in 32-bit binary IEEE 754 format (little
endian).

All strings are stored as a 32-bit integer string length (number of bytes)
followed by a UTF-8 format string (there is no zero termination and no BOM).

\section{File structure}
The structure of an OpenCTM file is as follows:

[Header]\newline
[Body data]

Each part of the file is described in the following chapters.


%-------------------------------------------------------------------------------

\chapter{Header}
The file must start with a header, which looks as follows:

\begin{tabular}{|l|l|l|}\hline
\textbf{Offset} &  \textbf{Type} & \textbf{Description}\\ \hline
0 & Integer & Magic identifier (0x4d54434f, or "OCTM" when read as ASCII).\\ \hline
4 & Integer & File format version (0x00000005 = version 5).\\ \hline
8 & Integer & Compression method, which can be one of the following:\\
 & & 0x00574152 - Use the RAW compression method.\\
 & & 0x0031474d - Use the MG1 compression method.\\
 & & 0x0032474d - Use the MG2 compression method.\\ \hline
12 & Integer & Vertex count.\\ \hline
16 & Integer & Triangle count.\\ \hline
20 & Integer & UV map count.\\ \hline
24 & Integer & Attribute map count.\\ \hline
28 & Integer & Boolean flags, or:ed together:\\
 & & 0x00000001 - The file contains per-vertex normals.\\ \hline
32 & String & File comment.\\ \hline
\end{tabular}

The length of the file header is varibale, since it contains fields of variable
length.


%-------------------------------------------------------------------------------

\chapter{Body data}
The body data follows immediately after the file header. Its file offset is
dictated by the length of the file header.

The format of the body data is specific for each compression method, which is
defined by the "Compression method" field in the header.

The body data contains the vertex, index, normal, UV map and attribute map
data, usually in a compressed form.


\section{RAW}
The layout of the body data for the RAW compression method is:

[Indices]\newline
[Vertices]\newline
[Normals]\newline
[UV Map 0]\newline
[UV Map 1]\newline
...\newline
[UV Map N]\newline
[Attribute Map 0]\newline
[Attribute Map 1]\newline
...\newline
[Attribute Map M]

\subsection{Indices}
The indices are stored as an integer identifier, 0x58444e49 ("INDX"), followed
by all the triangle indices. Each index is an unsigned integer value. There are
three indices per triangle, and the number of triangles is given by the
"Triangle count" field in the header:

\begin{tabular}{|l|l|l|}\hline
\textbf{Offset} &  \textbf{Type} & \textbf{Description}\\ \hline
0 & Integer & Identifier (0x58444e49, or "INDX" when read as ASCII).\\ \hline
4 & Integer & Vertex index for the 1st corner of the 1st triangle.\\ \hline
8 & Integer & Vertex index for the 2nd corner of the 1st triangle.\\ \hline
12 & Integer & Vertex index for the 3rd corner of the 1st triangle.\\ \hline
16 & Integer & Vertex index for the 1st corner of the 2nd triangle.\\ \hline
... & & \\ \hline
\end{tabular}

The length of the indices section is $4(1+3N)$ bytes, where $N$ is the triangle
count.

\subsection{Vertices}
The vertices are stored as an integer identifier, 0x54524556 ("VERT"), followed
by all the vertex coordinates. Each vertex coordinate is stored as three
floating point values ($x,y,z$), and the number of vertices is given by the
"Vertex count" field in the header:

\begin{tabular}{|l|l|l|}\hline
\textbf{Offset} &  \textbf{Type} & \textbf{Description}\\ \hline
0 & Integer & Identifier (0x54524556, or "VERT" when read as ASCII).\\ \hline
4 & Float & $x$ coordinate of the 1st vertex.\\ \hline
8 & Float & $y$ coordinate of the 1st vertex.\\ \hline
12 & Float & $z$ coordinate of the 1st vertex.\\ \hline
16 & Float & $x$ coordinate of the 2nd vertex.\\ \hline
... & & \\ \hline
\end{tabular}

The length of the vertices section is $4(1+3N)$ bytes, where $N$ is the vertex
count.

\subsection{Normals}
The normals section is optional, and only present if the per-vertex normals
flag is set in the header.

The normals are stored as an integer identifier, 0x4d524f4e ("NORM"), followed
by all the normal coordinates. Each normal is stored as three floating point
values ($x,y,z$), and the number of normals is given by the "Vertex count" field
in the header:

\begin{tabular}{|l|l|l|}\hline
\textbf{Offset} &  \textbf{Type} & \textbf{Description}\\ \hline
0 & Integer & Identifier (0x4d524f4e, or "NORM" when read as ASCII).\\ \hline
4 & Float & $x$ coordinate of the 1st normal.\\ \hline
8 & Float & $y$ coordinate of the 1st normal.\\ \hline
12 & Float & $z$ coordinate of the 1st normal.\\ \hline
16 & Float & $x$ coordinate of the 2nd normal.\\ \hline
... & & \\ \hline
\end{tabular}

The length of the normals section is $4(1+3N)$ bytes, where $N$ is the vertex
count.

\subsection{UV map}
To be written...

\subsection{Attribute map}
To be written...


\section{MG1}
To be written...

Please read the source code file compressMG1.c for more details on the body
format when the MG1 compression method is used.


\section{MG2}
To be written...

Please read the source code file compressMG2.c for more details on the body
format when the MG2 compression method is used.

\end{document}
